\documentclass[11pt]{article}

% =============================================================================
%  style.tex — the paper's design system.
%
%  Single-column modern-preprint layout: XCharter for text, Fira Sans for
%  headings and captions, Inconsolata for code. Colours are the same three
%  categorical slots the figures use (paper/figures/style.py), so a reader
%  never has to relearn what blue and orange mean between a figure and the
%  prose around it.
% =============================================================================

% ---------------------------------------------------------------- geometry --
\usepackage[
  letterpaper,
  top=1.05in, bottom=1.05in, left=1.10in, right=1.10in,
  footskip=0.45in, headsep=0.22in
]{geometry}

% ------------------------------------------------------------------- fonts --
\usepackage[T1]{fontenc}
\usepackage{XCharter}
\usepackage[scaled=0.93]{FiraSans}
\usepackage[varqu,varl,scaled=0.92]{inconsolata}
\usepackage[charter,vvarbb,scaled=1.04]{newtxmath}
\usepackage{microtype}

\linespread{1.045}
\setlength{\emergencystretch}{1.6em}   % long \texttt{} runs must not overflow

% ------------------------------------------------------------------ colours --
\usepackage[table,dvipsnames]{xcolor}

\definecolor{arblue}{HTML}{2A78D6}    % autoregressive  (categorical slot 1)
\definecolor{dgorange}{HTML}{EB6834}  % block diffusion (categorical slot 2)
\definecolor{llaqua}{HTML}{1BAF7A}    % LLaDA           (categorical slot 3)

\definecolor{ink}{HTML}{0B0B0B}
\definecolor{ink2}{HTML}{52514E}
\definecolor{muted}{HTML}{898781}
\definecolor{rule}{HTML}{C3C2B7}
\definecolor{hairline}{HTML}{E8E7E1}

\definecolor{washblue}{HTML}{EAF1FB}
\definecolor{washorange}{HTML}{FDEEE7}
\definecolor{washneutral}{HTML}{F4F3F0}

\definecolor{deepink}{HTML}{1B3A5C}    % heading ink — a calm dark blue

% ------------------------------------------------------------------ tables --
\usepackage{booktabs}
\usepackage{array}
\usepackage{multirow}
\usepackage{tabularx}
\renewcommand{\arraystretch}{1.18}
\setlength{\tabcolsep}{6pt}
\newcolumntype{R}{>{\raggedleft\arraybackslash}X}
\newcolumntype{L}{>{\raggedright\arraybackslash}X}

% Table body in the sans face at a slightly smaller size — numbers read
% better and the table separates cleanly from the surrounding text.
\newcommand{\tabfont}{\sffamily\footnotesize}
\newcommand{\thead}[1]{\textcolor{ink2}{\footnotesize\textbf{#1}}}

% --------------------------------------------------------------- headings ---
\usepackage{titlesec}

\titleformat{\section}
  {\normalfont\sffamily\large\bfseries\color{deepink}}
  {\textcolor{arblue}{\thesection}}{0.7em}{}
\titlespacing*{\section}{0pt}{1.55\baselineskip}{0.55\baselineskip}

\titleformat{\subsection}
  {\normalfont\sffamily\normalsize\bfseries\color{deepink}}
  {\textcolor{muted}{\thesubsection}}{0.6em}{}
\titlespacing*{\subsection}{0pt}{1.15\baselineskip}{0.35\baselineskip}

\titleformat{\subsubsection}
  {\normalfont\sffamily\small\bfseries\color{ink2}}
  {}{0em}{}
\titlespacing*{\subsubsection}{0pt}{0.9\baselineskip}{0.25\baselineskip}

\titleformat{\paragraph}[runin]
  {\normalfont\sffamily\small\bfseries\color{ink}}
  {}{0em}{}[.]
\titlespacing*{\paragraph}{0pt}{0.7\baselineskip}{0.55em}

% --------------------------------------------------------------- captions ---
\usepackage[font=small,labelfont={sf,bf,color=deepink},
            textfont={color=ink2},labelsep=period,
            skip=7pt,justification=justified]{caption}
\usepackage{subcaption}

% ------------------------------------------------------------------- boxes --
\usepackage[most]{tcolorbox}

% A key-result callout: tinted panel, coloured left rule, sans heading.
\newtcolorbox{keybox}[1][]{
  enhanced, breakable,
  colback=washneutral, colframe=washneutral,
  boxrule=0pt, leftrule=2.2pt,
  borderline west={2.2pt}{0pt}{arblue},
  arc=0pt, outer arc=0pt,
  left=9pt, right=9pt, top=7pt, bottom=7pt,
  before skip=9pt, after skip=9pt,
  fontupper=\small,
  #1
}

% A "what this section establishes" strip, used once per major section.
\newtcolorbox{takeaway}[1][]{
  enhanced, breakable,
  colback=washblue!55!white, colframe=washblue!55!white,
  boxrule=0pt, arc=1.5pt, outer arc=1.5pt,
  left=9pt, right=9pt, top=6pt, bottom=6pt,
  before skip=9pt, after skip=9pt,
  fontupper=\small,
  #1
}

% A caution/limitation strip — for scope conditions the reader must carry.
\newtcolorbox{caution}[1][]{
  enhanced, breakable,
  colback=washorange!60!white, colframe=washorange!60!white,
  boxrule=0pt, leftrule=2.2pt,
  borderline west={2.2pt}{0pt}{dgorange},
  arc=0pt, outer arc=0pt,
  left=9pt, right=9pt, top=6pt, bottom=6pt,
  before skip=9pt, after skip=9pt,
  fontupper=\small,
  #1
}

\newcommand{\boxhead}[1]{%
  \textcolor{deepink}{\sffamily\bfseries\footnotesize\MakeUppercase{#1}}\par
  \vspace{2pt}}

% ------------------------------------------------------------------- lists --
\usepackage{enumitem}
\setlist[itemize]{leftmargin=1.15em,itemsep=2.5pt,topsep=4pt,parsep=0pt}
\setlist[enumerate]{leftmargin=1.5em,itemsep=2.5pt,topsep=4pt,parsep=0pt}

% -------------------------------------------------------------------- misc --
\usepackage{amsmath}
\let\Bbbk\relax          % newtxmath already defines it
\usepackage{amssymb}
\usepackage{siunitx}
\sisetup{group-separator={,},group-minimum-digits=4,detect-all}
\usepackage{graphicx}
\usepackage{algorithm}
\usepackage{algpseudocode}
\usepackage{float}
\usepackage[hidelinks]{hyperref}
\hypersetup{
  colorlinks=true,
  linkcolor=arblue, citecolor=arblue, urlcolor=arblue,
  pdfborder={0 0 0}
}

% Model shorthands — one spelling per model, everywhere.
\newcommand{\ARthree}{\textsc{ar-3b}}
\newcommand{\ARoneb}{\textsc{ar-1.5b}}
\newcommand{\DG}{\textsc{dg}}
\newcommand{\LLaDA}{\textsc{llada-8b}}

% Coloured inline model tags, matching the figures.
\newcommand{\artag}{\textcolor{arblue}{\sffamily\bfseries AR}}
\newcommand{\dgtag}{\textcolor{dgorange}{\sffamily\bfseries DG}}

% A quiet inline source pointer.
\newcommand{\artifact}[1]{\textcolor{muted}{\ttfamily\scriptsize #1}}

% ---------------------------------------------------------------- title page --
\usepackage{titling}
\usepackage{fancyhdr}

\pagestyle{fancy}
\fancyhf{}
\renewcommand{\headrulewidth}{0pt}
\fancyfoot[C]{\textcolor{muted}{\sffamily\footnotesize\thepage}}

% \papertitle{title}{subtitle}{author}{affiliation}{version line}
\newcommand{\papertitle}[5]{%
  \thispagestyle{empty}
  \begingroup
  \raggedright\parindent=0pt
  {\sffamily\bfseries\fontsize{21}{25}\selectfont\color{deepink} #1\par}
  \vspace{5pt}
  {\sffamily\fontsize{13}{16}\selectfont\color{ink2} #2\par}
  \vspace{8pt}
  {\color{arblue}\rule{\linewidth}{2pt}\par}
  \vspace{9pt}
  {\sffamily\large\color{ink} #3\par}
  \vspace{2pt}
  {\sffamily\small\color{ink2} #4\par}
  \vspace{3pt}
  {\sffamily\footnotesize\color{muted} #5\par}
  \endgroup
  \vspace{16pt}
}

% paper/macros.tex
% LaTeX macros for numbers used across the paper and CHECKLIST.md.
% Shared so the paper and checklist cannot drift (#65 acceptance criterion).

% --- Pinned stack (#57) ---
\newcommand{\vllmVersion}{0.25.1}
\newcommand{\torchVersion}{2.11.0+cu130}
\newcommand{\gpuModel}{NVIDIA A100 80GB PCIe}

% --- Quality (#43 rescore, 9-field) ---
\newcommand{\arThreeBFnine}{0.8516}   % AR-3B at L=256, 9-field F1
\newcommand{\dgNineF}{0.8518}         % DG at L=256, 9-field F1
\newcommand{\arThreeBFtwelve}{0.8263}  % AR-3B at L=256, 12-field F1
\newcommand{\dgTwelveF}{0.7858}        % DG at L=256, 12-field F1

% --- Cost laws (Campaign B, pending re-measurement on pinned stack) ---
\newcommand{\arBase}{7.27}       % AR-3B fitted intercept (ms)
\newcommand{\arSlope}{6.29}      % AR-3B ms per output token
\newcommand{\dgBase}{378.8}      % DG base at K=8 (ms) -- was 418.24 at K=32
\newcommand{\dgStep}{119.6}     % DG per-canvas step (ms) -- #68 median of 7 step costs
\newcommand{\canvasLength}{256}   % DG canvas size (tokens)

% --- Crossovers (#64, calibrated on pinned stack) ---
\newcommand{\crossoverL}{59}      % crossover output length (tokens) at K=8
\newcommand{\crossoverLambda}{2.6} % crossover arrival rate (ev/s) at K=8

% --- Batching (#42, retitle) ---
\newcommand{\dgBatchMult}{1.18}  % DG batching multiplier (dInfer, cache='')

% --- Defect magnitudes (#65 / CHECKLIST.md) ---
\newcommand{\defectOneFBefore}{0.779}  % AR-1.5B F1 with repetition_penalty
\newcommand{\defectOneFAfter}{0.844}   % AR-1.5B F1 with penalty=1.0
\newcommand{\defectTwoBefore}{1841.6}  % AR-1.5B p50 HF loop (ms)
\newcommand{\defectTwoAfter}{413.3}    % AR-1.5B p50 vLLM (ms)
\newcommand{\defectThreeBefore}{0.599} % AR-3B L=512 last-wins F1
\newcommand{\defectThreeAfter}{0.826}  % AR-3B L=512 first-wins F1
\newcommand{\defectFourBefore}{0}      % AR-3B completeness (metric bug)
\newcommand{\defectFourAfter}{100}     % AR-3B completeness (fixed)

% --- Discrepancies (#57) ---
\newcommand{\arOldPfifty}{1617.6}  % AR-3B p50 Campaign B (ms)
\newcommand{\arNewPfifty}{1915.4}  % AR-3B p50 new stack (ms)
\newcommand{\dgOldBase}{418.24}    % DG base at K=32 (ms)
\newcommand{\dgNewBase}{378.8}     % DG base at K=8 (ms)

% --- Input-length slopes (#44) ---
\newcommand{\arInputSlope}{0.047}  % AR ms per prompt-token
\newcommand{\dgInputSlope}{0.123}   % DG ms per prompt-token (2.6x steeper)

% --- Base decomposition (#44 K-sweep) ---
\newcommand{\dgPrefillPct}{93.5}  % % of base that is prefill+fixed at K=8
\newcommand{\dgDenoisePct}{6.5}    % % of base that is denoising at K=8

\graphicspath{{../figures/}}

\begin{document}

\papertitle
  {Diffusion LLMs as Fast Models for Specific Tasks: The Cost Shape of Block-Structured Decoding}
  {When diffusion language models win on latency, and why they usually don't}
  {Aleksandr Kaigorodov}
  {Yandex}
  {Preprint \textperiodcentered\ August 2026 \textperiodcentered\ v2}

% ---------------------------------------------------------------- abstract --
\begingroup
\small
\noindent
{\sffamily\bfseries\footnotesize\color{deepink}ABSTRACT}\par
\vspace{3pt}
\noindent
Diffusion language models can be faster than autoregressive ones for
structured output --- but only in a narrow regime, and only after the
comparison is done fairly. We measure both architectures on a pinned serving
stack (vLLM, A100 80GB) for structured log extraction, where they reach
quality equivalence at 9-field macro F1. Their latency costs have different
\emph{shapes}: autoregressive cost is a line in output length with no floor;
block-diffusion cost is a staircase, flat inside a fixed-size block and
stepping at each boundary. The shapes cross at \crossoverL\ output tokens,
below which the autoregressive model is cheaper.

The diffusion model's advantage grows with output length: at $L{=}1024$ it is
$5.4\times$ faster at matched quality, uses less energy, and degrades under
deadline pressure by losing quality rather than completeness. But the advantage
does not survive deployment: at matched quality the autoregressive model is
\footprintRatio$\times$ smaller and converts concurrency into throughput far
better. What survives is one narrow regime --- long structured output under a
batch-1 constraint --- and we state the five conditions it requires.
\endgroup

\vspace{6pt}

% ================================================================ 1. intro ==
\section{Introduction}
\label{sec:intro}

Autoregressive (AR) language models generate text one token at a time; each
token requires a forward pass over the growing KV cache. Block-diffusion
models generate an entire block of tokens in parallel through iterative
denoising, then step to the next block. The structural difference produces
different cost shapes, and the shapes determine which architecture is faster
for a given workload.

We measure both on a pinned stack --- vLLM \vllmVersion{}, \gpuModel{}, BF16
--- for structured log extraction from the LOGHUB BGL corpus. The task is
12-field JSON extraction from log lines, scored on a 9-field macro F1. The
models are Qwen2.5-3B-Instruct (AR, 3.0\,B active) and
DiffusionGemma-26B-A4B (DG, 4.0\,B active of 26\,B total), matched on active
parameters. At $L{=}256$ output tokens they reach quality equivalence:
\arThreeBFnine{} vs \dgNineF{} on the 9-field metric ($\Delta = \qualityGap$).

Five defects in our own harness, four biased against the AR arm, were
corrected before any result was reported. All were found by re-reading
committed artifacts, not by running new experiments. The corrected numbers
are what follows.

% ========================================================== 2. cost laws ==
\section{The Two Cost Laws}
\label{sec:laws}

\subsection{Autoregressive: a line with no floor}

AR latency is linear in output length, with no floor:
\begin{equation}
\label{eq:ar}
T_{\textsc{ar}}(L_{\mathrm{out}}) = \arBase + \arSlope \cdot L_{\mathrm{out}}
\quad\text{(ms)}
\end{equation}
The intercept absorbs prompt prefill (amortised by prefix caching); the slope
is the per-token decode cost. Five repeats at $L{=}\canvasLength$ give p50
$=$ \arAnchorTwoFive\,ms (CV \arCV\%).

\subsection{Block diffusion: a staircase}

DiffusionGemma writes onto fixed-size canvases of \canvasLength{} tokens.
Within one canvas every position is denoised in parallel; crossing a boundary
triggers a fresh denoising pass. Cost is stepwise:
\begin{equation}
\label{eq:dg}
T_{\textsc{dg}}(L_{\mathrm{out}}) = \dgBase +
\dgStep \cdot \Bigl(\bigl\lceil \tfrac{L_{\mathrm{out}}}{\canvasLength}\bigr\rceil - 1\Bigr)
\quad\text{(ms)}
\end{equation}
The \dgBase\,ms base at $K{=}8$ denoising steps is mostly prefill (bounded at
53--68\% by two independent measurements; see Appendix~\ref{app:decomp}). Each
additional canvas costs \dgStep\,ms (median of 7 measured transitions, range
\dgStepLo--\dgStepHi). The staircase has sharp edges: growing output from 240
to \canvasLength{} tokens costs 0.2\,ms; growing from \canvasLength{} to 264
costs 121.7\,ms.

\begin{figure}[t]
\centering
\includegraphics[width=\linewidth]{figures/fig1_cost_shapes.pdf}
\caption{The two cost shapes. AR is a line; DG is a staircase. The shapes
cross at $L^{*} = \crossoverL$ output tokens.}
\label{fig:shapes}
\end{figure}

\subsection{The crossover}

At the reference prompt the crossing point is analytic:
\begin{equation}
L^{*} = \frac{\dgBase - \arBase}{\arSlope} = \crossoverL \text{ tokens},
\qquad
\lambda^{*} = \frac{1000}{\dgBase} = \crossoverLambda \text{ events/s}
\end{equation}
Below $L^{*}$ the AR model is cheaper per event; above it the diffusion model
is, until the next canvas edge. Both laws generalise from BGL to HDFS within
1.7\%.

% ========================================================== 3. deployment ==
\section{Where Diffusion Wins, and Where It Doesn't}
\label{sec:fair}

\subsection{The surviving advantage}

At forced output $L \geq 512$ tokens, the diffusion model is faster at
matched quality. Table~\ref{tab:niche} shows the gap. The advantage grows with
output length because the AR model pays per token while the diffusion model
pays per canvas: at $L{=}1024$ the gap is $5.4\times$.

\begin{table}[h]
\centering
\caption{Latency at matched quality (AR-1.5B vs DG-26B-A4B, batch-1).}
\label{tab:niche}
\begin{tabular}{@{}lrrr@{}}
\toprule
$L_{\mathrm{out}}$ & AR-1.5B (ms) & DG (ms) & DG advantage \\
\midrule
256  & \arOneFiveTwoFive{} & \dgTwoFive{}   & $2.7\times$ \\
512  & \arOneFiveFiveOne{} & \dgFiveOne{}   & $4.0\times$ \\
1024 & \arOneFiveTenTwoFour{} & \dgTenTwoFour{} & $5.4\times$ \\
\bottomrule
\end{tabular}
\end{table}

The diffusion model also degrades gracefully under deadline pressure: it
returns a complete canvas at lower quality, while the AR model returns a
truncated record that fails schema validation.

\subsection{Why the advantage usually doesn't survive}

Three factors close the niche:

\paragraph{Footprint.} The quality-matched AR model (Qwen2.5-1.5B, F1
\arOneFiveFnine) is \arOneFiveWeights\,GB; the diffusion model is
\dgWeights\,GB --- a \footprintRatio$\times$ difference.

\paragraph{Batching.} Continuous batching converts arrival concurrency into
throughput at \arBatchLo--\arBatchHi$\times$ for AR against
$\leq$\dgBatchVllmHi$\times$ for diffusion. Both diffusion measurements are
engine-bounded (vLLM caps the batch at 8; dInfer runs cache-less), so the
true ceiling is unknown but lies at or above \dgBatchVllmHi$\times$. Even so,
the gap is more than an order of magnitude (see Appendix~\ref{app:batching}).

\paragraph{Input length.} The diffusion prefill slope is
\inputSlopeRatio$\times$ steeper, because each denoising step re-attends the
full prompt. At longer inputs the base rises and $L^{*}$ moves upward.

\subsection{The five conditions}

The diffusion model is the right choice only when all five hold:
\begin{enumerate}[nosep]
  \item Output length $\geq 512$ tokens, forced.
  \item Batch-1 by requirement (tenant isolation, interactivity).
  \item Completeness is validated against a schema.
  \item The quality floor exceeds what the AR model delivers.
  \item A \dgWeights\,GB resident footprint is acceptable.
\end{enumerate}

% ========================================================== 4. related ==
\section{Related Work}
\label{sec:related}

Diffusion LM speedup claims (LLaDA~\cite{llada2025},
DiffusionGemma~\cite{diffusiongemma}) are typically measured against
author-built AR baselines. Constrained decoding (outlines~\cite{outlines2024},
xgrammar~\cite{xgrammar2024}, llguidance~\cite{llguidance}) is the standard
answer to the structured-output problem and is not measured here; it is a
limitation. Speculative decoding (Medusa~\cite{medusa2024},
EAGLE~\cite{eagle2024}) attacks the same per-token slope the AR cost law
turns on. Harness bias in systems benchmarking is discussed
by~\cite{hoefler2021}.

% ========================================================== 5. threats ==
\section{Threats to Validity}
\label{sec:threats}

All measurements are on one GPU (\gpuModel{}); the fitted constants are
A100-specific but the shapes should transfer. The batching ceiling reproduces
across two runtimes and two model families, but both measurements are
engine-bounded. The crossover is calibrated at a $\sim$\refPrompt-token prompt
and shifts with input length. The AR--DG quality difference ($\Delta =
\qualityGap$) has a 95\% CI of $[-0.012, +0.011]$ at $n=300$; the two are
statistically equivalent at a $\pm 0.01$ margin.

% ========================================================== 6. conclusion ==
\section{Conclusion}
\label{sec:conclusion}

Diffusion language models have a real latency advantage for long structured
output: $5.4\times$ faster at $L{=}1024$ at matched quality, with graceful
degradation under deadline pressure. The advantage comes from the cost shape
--- a staircase that is flat inside each block --- not from the diffusion
process itself. The floor is mostly prefill, and any decoder emitting
fixed-size blocks would inherit the same shape.

The advantage does not survive deployment. At matched quality the AR model is
\footprintRatio$\times$ smaller and batches more than an order of magnitude
better. What survives is a narrow regime: long forced output, batch-1, schema
validation, and a data-centre footprint. Outside that regime, the line beats
the staircase.

% ======================================================================
%                          APPENDICES
% ======================================================================

\appendix

% ------------------------------------------------------- A. harness defects
\section{Five Harness Defects}
\label{app:defects}

Each of the following was found by re-reading committed artifacts, not by
running new experiments. Four of the five biased against the autoregressive
arm; the fifth was a correctness-of-record defect.

\begin{enumerate}
\item \textbf{Decoding configuration.} Qwen2.5's checkpoint sets
\texttt{repetition\_penalty}\,$=$\,1.1, which suppresses span-copying and
empties the \texttt{host} field. Correcting to 1.0 moved AR 9-field F1 from
\defectOneFBefore{} to \defectOneFAfter. \emph{Against AR.}

\item \textbf{Serving parity.} The AR model ran a naive HuggingFace generate
loop while the diffusion model ran on vLLM with continuous batching. Moving
AR to vLLM moved its p50 from \defectTwoBefore{} to \defectTwoAfter\,ms
(\defectTwoRatio$\times$), reversing the headline speed claim.
\emph{Against AR.}

\item \textbf{Parser policy.} At $L{=}512$, \rambleRate\% of AR outputs
contained multiple JSON objects. The old parser took the last (most-drifted)
copy; first-object-wins moved AR F1 from \defectThreeBefore{} to
\defectThreeAfter. \emph{Against AR.}

\item \textbf{Metric defaults.} The completeness metric treated falsy values
(\texttt{""}, \texttt{0.0}, \texttt{False}) as missing keys. The AR model
emits falsy defaults on early stop; DiffusionGemma always fills every field.
Distinguishing absent from present-but-falsy moved AR completeness from
\defectFourBefore\% to \defectFourAfter\%. \emph{Against AR.}

\item \textbf{Provenance.} A CLI default wrote an AR model name into three DG
result files. Biases no comparison but made audit impossible until correction
sidecars were written. \emph{Neutral.}
\end{enumerate}

The clustering on one arm is not a coincidence: a comparison harness written by
someone hoping the new architecture wins fails asymmetrically, because the
arm the experimenter is not rooting for inherits the checkpoint defaults, the
naive serving path, and the mature production infrastructure.

% ------------------------------------------------------- B. decomposition
\section{Base Decomposition: Why the Fit Fails}
\label{app:decomp}

Sweeping $K \in \{2,4,8,16,32,64\}$ at fixed output length fits
$T(K) = \kfitBase + \kfitSlope \cdot K$ with $R^2 = \kfitRsq$. The low $R^2$ is
itself informative: latency saturates above $K{=}16$, so the linear term is a
poor description at the top of the range and the fit is carried by the bottom.
The intercept is therefore not interpretable as a floor.

\begin{figure}[h]
\centering
\includegraphics[width=0.85\linewidth]{figures/fig4_ksweep.pdf}
\caption{Denoising-step sweep. Latency saturates above $K{=}16$; quality is
flat above $K{=}8$.}
\label{fig:ksweep}
\end{figure}

Two measurement-based bounds do the work instead:

\begin{itemize}
\item \textbf{From $K{=}2$.} For any decomposition
$T(K) = \text{floor} + K \cdot c(K)$ with $c \geq 0$,
$\text{floor} \leq T(2) = 202$\,ms. At $K{=}8$ the denoising-attributable
share is $\geq 47\%$ of the \dgBase\,ms base.

\item \textbf{From the canvas step.} One additional canvas costs \dgStep\,ms
and is a $K{=}8$ denoising pass. So denoising is $\geq \dgStep / \dgBase = 32\%$
of the base. The first canvas's pass is somewhat cheaper than a marginal one,
but not by the $5\times$ the fitted intercept would require.
\end{itemize}

The honest range at $K{=}8$ is 32--47\% denoising against 53--68\% floor. The
exact split awaits profiler attribution.

% ------------------------------------------------------- C. batching
\section{The Batching Ceiling and Its Engine Defects}
\label{app:batching}

The diffusion batching ceiling ($\leq$\dgBatchVllmHi$\times$) reproduces across
two runtimes and two model families, but both measurements are bounded by
known engine defects:

\begin{itemize}
\item \textbf{vLLM} caps the batch at 8 when two requests mix denoising
phases, so the DiffusionGemma multipliers (\dgBatchVllmLo--\dgBatchVllmHi$\times$)
are lower bounds, not ceiling measurements.

\item \textbf{dInfer} (LLaDA-8B) errors on prompts of 632 tokens or more,
forcing a cache-less mode that inflates absolute latency by $\sim$30$\times$.
The multiplier (\dgBatchMult$\times$) is unaffected as a ratio, but the
claim that this is a cache-enabled measurement is not defensible.
\end{itemize}

\begin{figure}[h]
\centering
\includegraphics[width=\linewidth]{figures/fig6_batching.pdf}
\caption{Batch-1 to peak throughput multipliers. Both diffusion bars are
engine-bounded: vLLM caps the batch at 8, and dInfer runs cache-less.}
\label{fig:batching}
\end{figure}

The standing explanation --- that bidirectional attention invalidates KV-cache
reuse --- is partly falsified: disabling prefix caching on the AR arm drops
its multiplier from \cacheOnMult$\times$ to \cacheOffMult$\times$, still more
than twenty times the diffusion ceiling. Remaining candidates include
mixture-of-experts routing variability and kernel maturity. The true ceiling
on an uncrippled stack is unknown but lies at or above \dgBatchVllmHi$\times$.

\bibliographystyle{plain}
\bibliography{../references}

\end{document}
